\section{决策关联的价值编码}\label{chap3_7:results_choice_value}

\ref{chap3_6:results_offer_value}确立了财富对价值编码的调制作用。在此基础上，我们试图理解这一价值表征如何被进一步转换为最终选择，尤其关注其中介于选项价值与选择信号之间的中间表征。

价值差$\Delta V = V_E - V_S$是二元选择的自然连续模拟：$\Delta V > 0$时赚取选项占优，$\Delta V < 0$时消费选项占优，因此$\Delta V$是驱动选择的候选变量。我们首先考察$\Delta V$是否为单神经元水平上的主导编码特征。由于$\Delta V$是$V_E$和$V_S$的精确线性组合，线性回归存在固有的混淆。正则化虽能缓解共线性问题，但会倾向于用单一变量$\Delta V$进行解释，从而产生假阳性，无法确认其是否反映了神经元真实的编码模式。为此，我们采用类似图~\ref{fig5_2:joint}~的联合编码分析，对同时显著编码$V_E$和$V_S$的神经元进行考察：若神经元编码了$\Delta V$，则上述共编码神经元应不成比例地集中在反对角单元格（$V_{E}+ \times V_{S}-$或$V_{E-} \times V_{S}+$）。Fisher精确检验在选项呈现期和钱包更新期均未在任何脑区检测到显著的反对角集中模式（图~\ref{fig6_1:ve_vs}），说明$\Delta V$并非单神经元水平上的主导编码特征。

Padoa-Schioppa等人\citep{padoaschioppa2006neurons}在OFC发现了一类"被选价值（$V_{\mathrm{Chosen}}$）"神经元，其放电率与被选商品的主观价值成正比，而与具体选择了哪种商品无关。在赚取-消费范式中，编码$V_\mathrm{Chosen}$的神经元对$V_E$和$V_S$的调谐应依赖于选择方向：赚取试次中$V_E$主导，消费试次中$V_S$主导。图~\ref{fig6_2:example}A-D所示的OFC神经元呈现了这一模式的简单情形：$V_S$编码仅在消费试次中出现，赚取试次中消失，$V_E$则始终无显著调谐。图~\ref{fig6_2:example}E-H所示的另一个OFC神经元同时编码$V_E$和$V_S$，但方向和强度均依赖于选择——$V_S$在消费试次斜率为正、在赚取试次反转为负，$V_E$在赚取试次斜率更陡——与$V_\mathrm{Chosen}$编码的预期相符。我们将$V_E$、$V_S$、$V_\mathrm{Chosen}$、选择和财富水平同时纳入回归模型，统计有显著编码的神经元比例（图~\ref{fig6_3:proportion}）。结果表明，$V_\mathrm{Chosen}$在大多数脑区中的编码比例均超过随机水平。

\begin{figure}[H]
    \centering
    \includegraphics[width=0.6\textwidth]{Figure6_1_VE_VS}
    \bicaption{\enspace 赚取价值与消费价值的联合编码}{\enspace Joint encoding of earning and spending value}
    \label{fig6_1:ve_vs}
\end{figure}
\fignotetext{
    \textbf{(A) 选项呈现期。}四行依次为聚合数据、猕猴W、猕猴U、猕猴T；三列对应ACC、DLPFC、OFC。每个$2\times2$列联表展示同时对$V_E$和$V_S$有显著编码（$p < 0.05$）的神经元占该脑区所有神经元的比例，行为$V_E$编码方向（$V_E+$/$V_E-$），列为$V_S$编码方向（$V_S+$/$V_S-$）。对每个列联表进行双侧Fisher精确检验；若偏离程度接近无效（$p > 0.95$）则不绘制边框；否则根据优势方向绘制对角或反对角边框，偏离显著时边框为实线，不显著为虚线。所有面板均未出现反对角模式；凡检测到偏离独立性的面板均表现为对角集中（$V_E+/V_S+$与$V_E-/V_S-$共现）。\\
    \textbf{(B) 钱包更新期。}布局同(A)。所有面板均未出现反对角模式；凡检测到偏离独立性的面板均表现为对角集中（$V_E+/V_S+$与$V_E-/V_S-$共现）。}
    {\textbf{(A) Offer onset period.} The four rows show pooled data, Monkey W, Monkey U, and Monkey T; the three columns correspond to ACC, DLPFC, and OFC. Each $2\times2$ contingency table shows the proportion of neurons in each area with significant encoding of both $V_E$ and $V_S$ ($p < 0.05$), with rows indicating $V_E$ encoding direction ($V_E+$/$V_E-$) and columns indicating $V_S$ encoding direction ($V_S+$/$V_S-$). A two-sided Fisher exact test is applied to each table; if the association is statistically indistinguishable from independence ($p > 0.95$), no border is drawn; otherwise, a diagonal or anti-diagonal border is drawn according to the dominant pattern, with $p < 0.05$ shown as a solid border and $p \geq 0.05$ as a dashed border. No anti-diagonal pattern was observed in any table; all tables showing deviation from independence exhibited diagonal concentration ($V_E+/V_S+$ and $V_E-/V_S-$). \\
    \textbf{(B) Wallet update period.} Layout as in (A). No anti-diagonal pattern was observed in any table; all tables showing deviation from independence exhibited diagonal concentration ($V_E+/V_S+$ and $V_E-/V_S-$).}

\begin{figure}[H]
    \centering
    \includegraphics[width=0.9\textwidth]{Figure6_2_Example}
    \bicaption{\enspace 编码被选选项价值的示例神经元}{\enspace Representative neurons encoding chosen value}
    \label{fig6_2:example}
\end{figure}
\fignotetext{
    \textbf{(A, E) 选项呈现时刻的放电率。}试次按选择类型和$\Delta V$分组，赚取与消费试次内各自独立计算三分位线，共六组；实线对应赚取，虚线对应消费；阴影带为$\pm$标准误。颜色由红至绿按六组$\Delta V$均值升序排列，红色对应最高$\Delta V$，支持赚取选项，绿色对应最低$\Delta V$，支持消费选项。\\
    \textbf{(B, F) 与(A, E)相同分组方式下的调谐曲线。}纵轴为选项呈现后0–500~毫秒的平均放电率$\pm$标准误，横轴为各组$\Delta V$均值；同一选择类型内相邻分组以折线相连，颜色编码与(A, E)一致。\\
    \textbf{(C, G) 按选择类型和$V_E$三分位分组的调谐曲线。}在赚取与消费试次内各自独立计算$V_E$三分位线；纵轴与误差棒含义同(B, F)，横轴为各组$V_E$均值；颜色由浅红至深红按$V_E$均值升序排列。\\
    \textbf{(D, H) 按选择类型和$V_S$三分位分组的调谐曲线。}在赚取与消费试次内各自独立计算$V_S$三分位线；纵轴与误差棒含义同(B, F)，横轴为各组$V_S$均值；颜色由浅绿至深绿按$V_S$均值升序排列。}{
    \textbf{(A, E) Firing rate at offer onset.} Trials are grouped by choice type and $\Delta V$ tertile, computed separately within earning and spending trials, yielding six groups. Solid lines correspond to earning trials and dashed lines to spending trials; shading shows $\pm$ SEM. Color progresses from red (highest $\Delta V$) to green (lowest $\Delta V$) in ascending order of mean $\Delta V$ per group.\\
    \textbf{(B, F) Tuning curves using the same grouping as (A, E).} The y-axis shows mean firing rate $\pm$ SEM over 0–500~ms after offer onset; the x-axis shows mean $\Delta V$ per group. Adjacent groups within the same choice type are connected by lines; color coding as in (A, E).\\
    \textbf{(C, G) Tuning curves grouped by choice type and $V_E$ tertile.} $V_E$ tertiles are computed separately within earning and spending trials. Layout as in (B, F); x-axis shows mean $V_E$ per group. Color progresses from light red to dark red in ascending order of $V_E$.\\
    \textbf{(D, H) Tuning curves grouped by choice type and $V_S$ tertile.} $V_S$ tertiles are computed separately within earning and spending trials. Layout as in (C, G); x-axis shows mean $V_S$ per group. Color progresses from light green to dark green in ascending order of $V_S$.}

\begin{figure}[H]
    \centering
    \includegraphics[width=0.75\textwidth]{Figure6_3_Proportion}
    \bicaption{\enspace 显著编码各任务变量的神经元比例}{\enspace Proportion of neurons significantly encoding task variables}
    \label{fig6_3:proportion}
\end{figure}
\fignotetext{四行依次为聚合数据、猕猴W、猕猴U、猕猴T。以选项呈现后0–500~毫秒内的平均放电率为因变量做线性回归。彩色柱为显著编码神经元的比例（$\alpha = 0.05$）；灰色柱为打乱试次标签所得的经验假阳性率。红色、绿色、蓝色分别对应ACC、DLPFC和OFC。}{Four rows show data for pooled, Monkey W, Monkey U, and Monkey T. Mean firing rate over 0–500~ms after offer onset was used. Colored bars show the proportion of neurons with significant encoding ($\alpha = 0.05$); gray bars show the empirical false-positive rate estimated by permuting trial labels. Red, green, and blue correspond to ACC, DLPFC, and OFC, respectively.}

在单神经元分析的基础上，我们进一步在群体水平检验价值表征是否依赖选择类型，我们分别在赚取和消费试次中独立拟合TDR价值编码轴（图~\ref{fig6_4:conditioned}）。结果显示，“一致”条件（编码变量与选择类型匹配：赚取试次的$V_E$轴、消费试次的$V_S$轴）的编码强度显著高于"不一致"条件（编码变量与选择类型不匹配），表明前额叶群体的价值编码是选择依赖的，与单神经元水平的分析结果一致。

\begin{figure}[H]
    \centering
    \includegraphics[width=0.95\textwidth]{Figure6_4_ConditionedTDR}
    \bicaption{\enspace 选择条件化的TDR分析}{\enspace Choice-conditioned TDR analysis}
    \label{fig6_4:conditioned}
\end{figure}
\fignotetext{$2\times2$布局：左上为聚合数据，右上为猕猴W，左下为猕猴U，右下为猕猴T。每组内第一行为以$R^2$衡量$V_E$编码强度（均值$\pm$标准差），第二行$V_S$；第一列对应赚取试次，第二列对应消费试次。TDR分别在赚取和消费试次池中独立拟合。聚合数据的纵轴固定为$[0, 1]$，个体数据在每组内统一纵轴范围。一致条件（赚取试次的$V_E$轴、消费试次的$V_S$轴，对角线位置）的$R^2$高于不一致条件（反对角线位置）在聚合数据和猕猴W、U中均可见；猕猴T的$V_E$编码整体较弱，赚取和消费试次间无明显差异，但$V_S$的选择依赖模式与其他数据一致。}{$2\times2$ layout: top-left shows pooled data, top-right Monkey W, bottom-left Monkey U, bottom-right Monkey T. Within each panel, row 1 shows encoding strength measured by $R^2$ for $V_E$ (mean $\pm$ SD) and row 2 for $V_S$; column 1 corresponds to earn trials and column 2 to spend trials. TDR axes were fitted independently within earn and spend trial pools. The y-axis is fixed at $[0, 1]$ for pooled data and unified within each individual's panel. $R^2$ is higher for congruent conditions (the $V_E$ axis in earn trials and the $V_S$ axis in spend trials, diagonal positions) than for incongruent conditions (off-diagonal positions), visible in pooled data and Monkeys W and U. Monkey T shows weak $V_E$ encoding overall with no clear earn/spend difference, but the choice-dependent pattern for $V_S$ is consistent with the other panels.}

$V_\mathrm{Chosen}$的编码可能作为驱动选择的前向信号在选择发生之前形成，也可能只是选择完成后的下游反馈。要区分这两种可能，最直接的方式是比较价值信号与选择信号出现的时间顺序。我们对TDR得到的每条$R^2(t)$曲线分别计算信号潜伏期，图~\ref{fig6_5:latency}A展示了$t_{onset}$分布及各变量间的统计比较结果。在聚合数据中，所有三个脑区的$V_E$、$V_S$和$V_\mathrm{Chosen}$的$t_{onset}$均不晚于选择信号。由于信号潜伏期对噪声较为敏感，我们同时计算半峰时间（$t_{50}$）作为更稳健的对照指标（图~\ref{fig6_5:latency}A），结果与$t_{onset}$一致，说明时间顺序的结论不依赖于具体指标的选择。这一时间顺序表明$V_\mathrm{Chosen}$信号更有可能参与决策形成而非仅仅反映决策结果。

\begin{figure}[H]
    \centering
    \includegraphics[width=\textwidth]{Figure6_5_Latency}
    \bicaption{\enspace 价值信号与选择信号的涌现时间比较}{\enspace Temporal ordering of value and choice signals}
    \label{fig6_5:latency}
\end{figure}
\fignotetext{
    \textbf{(A) 信号潜伏期（$t_{onset}$）。}三行分别对应ACC、DLPFC、OFC；四列从左至右依次为聚合数据、猕猴W、猕猴U、猕猴T。信号潜伏期定义为选项呈现后$R^2$首次持续连续$\geq50$~毫秒超过刺激前基线阈值（第95百分位数）的时刻。每个子图以水平小提琴图展示四个变量（$V_E$、$V_S$、$V_\mathrm{Chosen}$、选择）在100次交叉验证下的分布；小提琴面积正比于有效数据次数——无持续显著时段时为无效；横线标示中位数和四分位距。各价值变量与选择信号的比较采用双侧Wilcoxon秩和检验（$\alpha = 0.05$），$p$值在价值信号更早时显示在左侧，选择信号更早时显示在右侧，无显著差异时居中。信号潜伏期对噪声较为敏感，部分子图中有效置换比例较低、分布较宽。\\
    \textbf{(B) 半峰潜伏期（$t_{50}$）。}布局同(A)。$t_{50}$定义为$R^2$首次超过半峰值的时刻（基线中位数与峰值的均值），峰值未达显著时为无效。统计比较方法同(A)。相较于(A)，$t_{50}$对噪声更不敏感，但两者反映的时间顺序总体一致。}
    {\textbf{(A) Signal latency.} The three rows correspond to ACC, DLPFC, and OFC; the four columns show pooled data, Monkey W, Monkey U, and Monkey T. Onset latency is defined as the first time point at which $R^2$ continuously exceeds the 95th percentile of the pre-stimulus baseline for $\geq 50$~ms. Each panel shows horizontal violin plots of the four variables ($V_E$, $V_S$, $V_\mathrm{Chosen}$, choice) across 100 cross validations; violin area is proportional to the fraction of permutations yielding a valid estimate — onset latency is invalid when no sustained significant period exists. Horizontal lines indicate the median and interquartile range. Pairwise comparisons between each value variable and the choice signal used two-sided Wilcoxon rank-sum tests ($\alpha = 0.05$); $p$-values are displayed on the left when the value signal is earlier, on the right when the choice signal is earlier, and centered when not significant. Onset latency is more sensitive to noise; some panels show lower valid fractions and wider distributions. \\
    \textbf{(B) Time to half-peak ($t_{50}$).} Layout as in (A). $t_{50}$ is defined as the first time $R^2$ exceeds the half-peak value (mean of baseline median and peak), computed only when the peak is significant; $t_{50}$ is invalid when the peak is not significant. Statistical comparisons followed the same procedure as in (A). Compared with $t_{onset}$, $t_{50}$ is less sensitive to noise and the temporal ordering is consistent between the two metrics.}

最后，我们比较了三个脑区对$V_E$、$V_S$和$V_\mathrm{Chosen}$的编码强度（图~\ref{fig6_6:comparison}）。整体而言，编码强度按$V_S > V_\mathrm{Chosen} > V_E$递减，但三种变量各自呈现出不同的脑区分布模式。$V_E$的编码最弱，所有个体的ACC中均未发现$V_E$编码，而OFC对$V_E$的编码强度在所有个体中均高于DLPFC和ACC；ACC与OFC的$V_S$编码强度大致持平，而DLPFC明显偏低；$V_\mathrm{Chosen}$的编码没有表现出明显的脑区差异。值得注意的是，第~\ref{chap3_3:results_choice}~节已表明OFC对物品域选择信号也有较强的编码，综合来看，OFC是三个脑区中唯一同时可靠编码$V_E$、$V_S$、$V_\mathrm{Chosen}$和选择信号的脑区，提示其可能是从价值计算到行为选择转化的关键脑区。

综合来看，前额叶皮层在单神经元水平并未编码$\Delta V$，但是以选择依赖的方式编码被选选项的价值$V_\mathrm{Chosen}$。这种表征在时间上先于选择信号的出现，暗示该价值编码可能参与了选择的形成，而非仅反映其结果。其中OFC最为广泛地编码了多种价值信号以及物品域选择结果，可能在价值到选择的转化过程中发挥了核心作用。

\begin{figure}[H]
    \centering
    \includegraphics[width=\textwidth]{Figure6_6_Comparison}
    \bicaption{\enspace 各脑区对价值变量的编码强度比较}{\enspace Encoding strength of value variables across brain areas}
    \label{fig6_6:comparison}
\end{figure}
\fignotetext{$4$行$\times$3列：四行分别为聚合数据、猕猴W、猕猴U和猕猴T，三列分别对应$V_E$、$V_S$和$V_\mathrm{Chosen}$。每个子图在$[-200, 500]$~毫秒的时间窗内叠加展示三个脑区的$R^2(t)$（均值$\pm$标准差），ACC为红色，DLPFC为绿色，OFC为蓝色。聚合数据的纵轴固定为$[0, 1]$；个体数据在各行内统一纵轴范围，以便比较同一被试内三种价值变量的编码强度。}{$4\text{-row} \times 3\text{-column}$ layout: four rows show pooled data, Monkey W, Monkey U, and Monkey T; three columns correspond to $V_E$, $V_S$, and $V_\mathrm{Chosen}$. Each panel shows $R^2(t)$ (mean $\pm$ SD) for the three brain areas over $[-200, 500]$~ms, with ACC in red, DLPFC in green, and OFC in blue. The y-axis range is fixed at $[0, 1]$ for pooled data; within each individual's row, the y-axis range is unified across the three value variables to facilitate comparison of relative encoding strength within the same monkey.}
